\clearpage
\fancyhead[L]{\small GERLING. SONNENFINSTERNISS.}
\hypertarget{native-627-lower}{}
\begin{center}\small Göttingische gelehrte Anzeigen. Stück 103. Seite 1025..1030. 1820 Juni 26.\\[3pt]\rule{.72\linewidth}{.3pt}\end{center}

Göttingen. Ueber die am 7. Sept. d. J. einfallende grosse Sonnenfinsterniss hat Hr. Prof. \textsc{Gerling} in Marburg, welcher dieselbe bereits vor acht Jahren in seiner Inaugural-Dissertation behandelt hatte (s. unsre Anz. 1812 S. 1889 [S. 542 d. B.]), in einem Schreiben an den Hrn. Hofr. \textsc{Gauss}, vom 23. Mai, die Resultate seiner wiederholten und weiter ausgeführten Untersuchungen mitgetheilt. Da diese Finsterniss im grössten Theile von Deutschland ringförmig erscheinen und für diese Gegenden, auf lange Zeit, das merkwürdigste Phänomen dieser Art sein wird, so glauben wir vielen Lesern dieser Blätter \SourceLeaf{628} durch die Bekanntmachung eines Auszuges aus jenen Mittheilungen einen willkommenen Dienst zu erweisen.

Folgende aus den \textsc{Carlini}'schen Sonnentafeln und den \textsc{Bürg}'schen Mondtafeln entlehnten Sonnen- und Mondsörter dienen der Rechnung zur Grundlage:

1820 September 7. Pariser mittl. Zeit.

\begin{center}
\small\setlength{\tabcolsep}{2pt}\begin{tabular}{l|rrr|rrr|rrr}
 & \multicolumn{3}{c|}{\shortstack{Länge der\\Sonne.}} & \multicolumn{3}{c|}{\shortstack{Länge des\\Mondes.}} & \multicolumn{3}{c}{\shortstack{Nordl. Breite\\des Mondes.}} \\
\hline
$11^{\mathrm u}$ Vorm. & $164^\circ$ & $40'$ & $26''4$ & $163^\circ$ & $19'$ & $53''6$ & $0^\circ$ & $52'$ & $45''0$ \\
2 Nachm. & & $47$ & $44.0$ & $164$ & $48$ & $17.7$ & & $44$ & $41.6$ \\
5 & & $55$ & $1.7$ & $166$ & $16$ & $38.3$ & & $36$ & $34.9$ \\
\end{tabular}
\end{center}

\begin{center}\small\setlength{\tabcolsep}{3pt}
\begin{tabular}{lr}
Horizontalparallaxe des Mondes . . . . . . . . . . & $53'\ 55''5$ \\
\hspace{1em}---\hspace{1em}der Sonne . . . . . . . . . . & $8.7$ \\
Halbmesser des Mondes . . . . . . . . . . . . . & $14'\ 43.1$ \\
\hspace{1em}---\hspace{1em}der Sonne . . . . . . . . . . . . . & $15\ 54.8$ \\
Schiefe der Ekliptik . . . . . . . . . . . . . & $23^\circ\ 27\ 56.2$ \\
\end{tabular}
\begin{tabular}{ll}
Die Conjunction erfolgt demnach um . . . . . . & $1^{\mathrm u}\ 58^{\mathrm m}\ 46^{\mathrm s}3$ M. Zeit in Paris \\
in der Länge . . . . . . . . . . . . . . . & $164^\circ\ 47'\ 41''$ \\
\end{tabular}
\end{center}

Für folgende einzelne Oerter hat Hr. Prof. \textsc{Gerling} zur Erleichterung der Beobachtungen die Hauptmomente der Erscheinung im Voraus berechnet:

\begin{center}\small\setlength{\tabcolsep}{3pt}
\begin{tabular}{l|rrr|rrr|rrr}
 &  \multicolumn{3}{c|}{Anfang.}  &  \multicolumn{3}{c|}{Mittel.}  &  \multicolumn{3}{c|}{Ende.}  \\
\hline
Berlin  &  $1^{\mathrm u}$  &  $31^{\mathrm m}$  &  $55^{\mathrm s}$  &  $2^{\mathrm u}$  &  $57^{\mathrm m}$  &  $31^{\mathrm s}$  &  $4^{\mathrm u}$  &  $15^{\mathrm m}$  &  $39^{\mathrm s}$  \\
Bremen  &  $1$  &  $7$  &  $13$  &  $2$  &  $33$  &  $38$  &  $3$  &  $54$  &  $5$  \\
Breslau  &  $1$  &  $53$  &  $36$  &  $3$  &  $18$  &  $29$  &  $4$  &  $34$  &  $54$  \\
Cassel  &  $1$  &  $13$  &  $39$  &  $2$  &  $40$  &  $30$  &  $4$  &  $0$  &  $54$  \\
Cöln  &  $1$  &  $1$  &  $6$  &  $2$  &  $28$  &  $52$  &  $3$  &  $50$  &  $16$  \\
Danzig  &  $1$  &  $55$  &  $32$  &  $3$  &  $19$  &  $5$  &  $4$  &  $34$  &  $22$  \\
Göttingen  &  $1$  &  $15$  &  $51$  &  $2$  &  $42$  &  $28$  &  $4$  &  $2$  &  $39$  \\
Gotha  &  $1$  &  $20$  &  $59$  &  $2$  &  $47$  &  $34$  &  $4$  &  $7$  &  $30$  \\
Halle  &  $1$  &  $26$  &  $26$  &  $2$  &  $52$  &  $20$  &  $4$  &  $11$  &  $45$  \\
Hamburg  &  $1$  &  $12$  &  $26$  &  $2$  &  $38$  &  $19$  &  $3$  &  $58$  &  $14$  \\
\end{tabular}
\end{center}
\begin{center}\small\setlength{\tabcolsep}{3pt}
\begin{tabular}{l|rrr|rrr|rrr}
  &  \multicolumn{3}{c|}{Anfang.}  &  \multicolumn{3}{c|}{Mittel.}  &  \multicolumn{3}{c}{Ende.}  \\
\hline
 Hannover  &  $1^{\mathrm u}$  &  $13^{\mathrm m}$  &  $11^{\mathrm s}$  &  $2^{\mathrm u}$  &  $39^{\mathrm m}$  &  $36^{\mathrm s}$  &  $3^{\mathrm u}$  &  $59^{\mathrm m}$  &  $45^{\mathrm s}$ \\
 Königsberg  &  $2$  &  $4$  &  $43$  &  $3$  &  $27$  &  $20$  &  $4$  &  $41$  &  $32$ \\
 Leipzig  &  $1$  &  $28$  &  $38$  &  $2$  &  $54$  &  $25$  &  $4$  &  $13$  &  $45$ \\
 Mannheim  &  $1$  &  $11$  &  $47$  &  $2$  &  $39$  &  $36$  &  $4$  &  $0$  &  $31$ \\
 Marburg  &  $1$  &  $11$  &  $23$  &  $2$  &  $38$  &  $40$  &  $3$  &  $59$  &  $18$ \\
 München  &  $1$  &  $30$  &  $40$  &  $2$  &  $57$  &  $47$  &  $4$  &  $17$  &  $33$ \\
 Ofen  &  $2$  &  $11$  &  $2$  &  $3$  &  $35$  &  $49$  &  $4$  &  $51$  &  $31$ \\
 Prag  &  $1$  &  $41$  &  $51$  &  $3$  &  $7$  &  $50$  &  $4$  &  $25$  &  $46$ \\
 Tübingen  &  $1$  &  $16$  &  $42$  &  $2$  &  $44$  &  $36$  &  $4$  &  $5$  &  $24$ \\
 Wien  &  $1$  &  $55$  &  $45$  &  $3$  &  $21$  &  $24$  &  $4$  &  $38$  &  $28$ \\
\end{tabular}
\end{center}

Die Zeiten sind hier und in der folgenden Tafel in \emph{wahrer Sonnenzeit} angesetzt.

\begin{center}\small\setlength{\tabcolsep}{3pt}
\begin{tabular}{l|r|r|rrr|rrr}
 &  $P$  &  $D$  &  \multicolumn{3}{c|}{\shortstack{Erste innere\\Berührung.}}  &  \multicolumn{3}{c|}{\shortstack{Zweite innere\\Berührung.}}  \\
\hline
Berlin  &  $72^\circ$  &  $84''$  &  ---  &  ---  &   &  ---  &  ---  &   \\
Bremen  &  $65$  &  $9$  &  $2^{\mathrm u}$  &  $30'$  &  $46''$  &  $2^{\mathrm u}$  &  $36^{\mathrm m}$  &  $30^{\mathrm s}$  \\
Breslau  &  $78$  &  $131$  &  ---  &  ---  &   &  ---  &  ---  &   \\
Cassel  &  $66$  &  $1$  &  $2$  &  $37$  &  $37$  &  $2$  &  $43$  &  $24$  \\
Cöln  &  $62$  &  $55$  &  $2$  &  $27$  &  $23$  &  $2$  &  $30$  &  $21$  \\
Danzig  &  $79$  &  $186$  &  ---  &  ---  &   &  ---  &  ---  &   \\
Göttingen  &  $67$  &  $10$  &  $2$  &  $39$  &  $37$  &  $2$  &  $45$  &  $19$  \\
Gotha  &  $68$  &  $18$  &  $2$  &  $44$  &  $47$  &  $2$  &  $50$  &  $20$  \\
Halle  &  $70$  &  $48$  &  $2$  &  $50$  &  $20$  &  $2$  &  $54$  &  $19$  \\
Hamburg  &  $66$  &  $36$  &  $2$  &  $35$  &  $55$  &  $2$  &  $40$  &  $43$  \\
\end{tabular}
\end{center}
\begin{center}\small\setlength{\tabcolsep}{3pt}
\begin{tabular}{l|r|r|rrr|rrr}
  &  $P$  &  $D$  &  \multicolumn{3}{c|}{\shortstack{Erste innere\\Berührung.}}  &  \multicolumn{3}{c}{\shortstack{Zweite innere\\Berührung.}}  \\
\hline
 Hannover  &  $66^\circ$  &  $17''$  &  $2^{\mathrm u}$  &  $36^{\mathrm m}$  &  $49^{\mathrm s}$  &  $2^{\mathrm u}$  &  $42^{\mathrm m}$  &  $23^{\mathrm s}$ \\
 Königsberg  &  $81$  &  $217$  &  ---  &  ---  &   &  & --- &  \\
 Leipzig  &  $71$  &  $52$  &  $2$  &  $52$  &  $49$  &  $2$  &  $56$  &  $1$ \\
 Mannheim  &  $64$  &  $44$  &  $2$  &  $37$  &  $28$  &  $2$  &  $41$  &  $43$ \\
 Marburg  &  $65$  &  $18$  &  $2$  &  $35$  &  $53$  &  $2$  &  $41$  &  $27$ \\
 München  &  $71$  &  $2$  &  $2$  &  $54$  &  $54$  &  $3$  &  $0$  &  $40$ \\
 Ofen  &  $84$  &  $123$  &  ---  &  ---  &   &  ---  &  ---  &  \\
 Prag  &  $75$  &  $75$  &  ---  &  ---  &   &  ---  &  ---  &  \\
 Tübingen  &  $66$  &  $46$  &  $2$  &  $42$  &  $35$  &  $2$  &  $46$  &  $37$ \\
 Wien  &  $79$  &  $87$  &  ---  &  ---  &   &  ---  &  ---  &  \\
\end{tabular}
\end{center}

Es bezeichnet hier $P$ den Punkt des Sonnenrandes, wo beim Anfang der Finsterniss der Mond zuerst eingreift, indem man die Grade vom nördlichsten Punkte (in Beziehung auf den Verticalkreis) nach \SourceLeaf{629} Westen zu zählt; $D$ hingegen gibt in Bogensecunden die scheinbare Entfernung der Mittelpunkte der Sonne und des Mondes, für das Mittel der Finsterniss, an. Die Abplattung ist bei allen Rechnungen zu $\tfrac{1}{300}$ angenommen.

Die Hauptmomente der Finsterniss für andere Oerter in Deutschland können näherungsweise aus denjenigen, welche in obiger ersten Tafel für einen zunächst liegenden angesetzt sind, vermittelst folgender Formeln abgeleitet werden, wo d$L$ den Längenunterschied beider Oerter in Zeitsecunden, d$B$ den Breitenunterschied in Bogensecunden und d$Z$ die Aenderung bedeutet, die den Angaben der Tafel beigefügt werden muss. Hiebei ist d$L$ positiv zu nehmen, wenn der Ort, für welchen man die Bestimmung wünscht, östlich von dem gewählten Ort der Tafel liegt, und d$B$ gleichfalls positiv, wenn jener nördlich von diesem ist.

Für den Anfang
\[
\mathrm{d}Z = +1.313\,\mathrm{d}L - 0.0296\,\mathrm{d}B
\]
Für das Mittel
\[
\mathrm{d}Z = +1.242\,\mathrm{d}L - 0.0343\,\mathrm{d}B
\]
Für das Ende
\[
\mathrm{d}Z = +1.118\,\mathrm{d}L - 0.0363\,\mathrm{d}B
\]

Besonders wichtig ist die Kenntniss der Grenzlinien der Zone, innerhalb welcher die Finsterniss ringförmig erscheint. Hr. Prof. \textsc{Gerling} hat folgende Punkte in diesen Grenzlinien durch Rechnung bestimmt, wo die Längen von Ferro an gerechnet sind.

\begin{center}
\begin{tabular}{c|cc|cc}
\shortstack{Wahre\\Zeit in\\Paris.} & \multicolumn{2}{c|}{\shortstack{Berührung der\\südl. Ränder.}} & \multicolumn{2}{c}{\shortstack{Berührung der\\nördl. Ränder.}} \\
 & Breite. & Länge. & Breite. & Länge. \\
\hline
$2^{\mathrm u}\ 4^{\mathrm m}$ & & & $52^\circ\ 59'$ & $22^\circ\ 43'$ \\
6 & & & $52\ \ 15$ & $23\ \ 15$ \\
8 & $54^\circ\ 12'$ & $28^\circ\ 50'$ & $51\ \ 31$ & $23\ \ 46$ \\
10 & $53\ \ 26$ & $29\ \ 18$ & $50\ \ 47$ & $24\ \ 17$ \\
12 & $52\ \ 40$ & $29\ \ 47$ & $50\ \ \ 4$ & $24\ \ 48$ \\
14 & $51\ \ 55$ & $30\ \ 16$ & $49\ \ 21$ & $25\ \ 19$ \\
16 & $51\ \ 10$ & $30\ \ 46$ & $48\ \ 39$ & $25\ \ 49$ \\
18 & $50\ \ 26$ & $31\ \ 15$ & $47\ \ 57$ & $26\ \ 19$ \\
20 & $49\ \ 42$ & $31\ \ 45$ & $47\ \ 15$ & $26\ \ 49$ \\
22 & $48\ \ 58$ & $32\ \ 15$ & $46\ \ 34$ & $27\ \ 20$ \\
24 & $48\ \ 15$ & $32\ \ 46$ & $45\ \ 53$ & $27\ \ 50$ \\
26 & $47\ \ 32$ & $33\ \ 17$ & & \\
28 & $46\ \ 50$ & $33\ \ 49$ & & \\
\end{tabular}
\end{center}

Verbindet man diese Punkte auf einer Landkarte durch gerade Linien, so trifft man folgende Oerter, von denen keiner mehr als etwa $2\tfrac{1}{2}$ Meilen von den Grenzlinien entfernt liegt. Die durch den Druck ausgezeichneten Oerter liegen der Grenze am nächsten.

Für die Berührung der \emph{südlichen} Ränder: Heiligenhafen, Neustadt, Grevesmühlen, Wismar, Schwerin, Parchim, Grabow, \emph{Perleberg}, Werben, Havelberg, \emph{Arneburg}, Tangermünde, Rathenow, Genthin, Ziesar, Coswig, Wittenberg, \emph{Kemberg}, Düben, Torgau, Wurzen, Oschatz, Döbeln, \emph{Freiberg}, Seyda, Kommotau, \emph{Rakonitz}, Przibram, \emph{Teyn}, Budweis, \emph{Grazen}, \emph{Ips}, \emph{Oberndorf}, Tirnitz, Fürstenfeld u.\,s.\,w.

Für die Berührung der \emph{nördlichen} Ränder: Staveren, Medenblick, Harderwyk, Amersfort, \emph{Barneveld}, Arnheim, \emph{Nimwegen}, Cleve, Grave, \emph{Genep}, Geldern, \emph{Venloo}, Crefeld, Roermunde, \emph{Erkelens}, Jülich, \emph{Nidegen}, Münster-Eiffel, Aldenau, Mondheim, \emph{Zell}, Trarbach, Kirchberg, \emph{Oberstein}, Münchweiler, \SourceLeaf{630} Kaiserslautern, Bergzabern, Hagenau, \emph{Stollhofen}, Achern, Freydenstadt, Hornberg, Rottweil, Villingen, \emph{Geisingen}, Schaffhausen, \emph{Stein}, \emph{Frauenfeld}, Winterthur, Chur u.\,s.\,w.

Es ist sehr zu wünschen, dass in der Nähe der Grenzlinien recht viele Beobachtungen angestellt werden, die schon dadurch einen Werth erhalten, dass unter Angabe des Beobachtungsplatzes bemerkt wird, ob sich ein Ring gebildet habe oder nicht, und im ersten Falle, wie lange die Erscheinung des Ringes gedauert habe. Da hier nur von Abmessung eines sehr kurzen Zeitraums die Rede ist, so kann allenfalls dabei der Mangel einer Secundenuhr durch das Abzählen der Schläge einer gewöhnlichen Taschenuhr oder der Schwingungen einer an einem Faden aufgehängten schweren Kugel ersetzt werden, wenn nur im ersten Fall die Anzahl der auf eine Minute gehenden Schläge, und im zweiten die Länge des Fadens und die Grösse der Kugel in irgend einem bekannten Masse mit angezeigt werden. Mit einem guten Fernrohr müsste aber doch eine solche Beobachtung angestellt sein, wenn sie einen Werth haben soll.

Da auch an den Orten, wo die Finsterniss ringförmig erscheint, noch völlig ein Achtel der Sonnenscheibe unverfinstert bleibt, so ist eine starke Abnahme des Tageslichts nicht zu erwarten, noch weniger, dass Sterne dem unbewaffneten Auge sichtbar werden sollten. Zum Besten solcher Personen, die während der grössten Verfinsterung einige der der Sonne nächsten und hellsten Fixsterne mit Fernröhren aufzusuchen oder mit guten Reflexionswerkzeugen ihre Distanzen vom Mondsrande zu messen versuchen möchten, hat Hr. Prof. \textsc{Gerling} noch folgende Angaben beigefügt, die zur Erleichterung dieser Aufsuchung dienen können.

\begin{center}
\begin{tabular}{lcc}
Für Regulus & $P = 65\tfrac{3}{4}^\circ$ Westl., & $D = 17\tfrac{1}{2}^\circ$ \\
Für $\ObservedBeta$ Löwe & $P = 41\tfrac{3}{4}^\circ$ Oestl., & $D = 13^\circ$ \\
Für Spica & $P = 116^\circ$ Oestl., & $D = 36\tfrac{1}{2}^\circ$. \\
\end{tabular}
\end{center}

Es bedeutet hier $D$ die Entfernung des Sterns vom Mittelpunkt der Sonne, und $P$ den Winkel, welchen der grösste Kreis vom Sonnenmittelpunkt zum Stern mit dem grössten Kreise vom Sonnenmittelpunkt zum Nordpol macht. Die Angaben sind für $2^{\mathrm u}\ 15^{\mathrm m}$ W.\,Z. in Paris berechnet, können aber in ganz Deutschland für die Zeit der grössten Verfinsterung dienen.
